% Lösungen Geraden (zusammenbau v0.4, Heft)
\einheitenkopf{Geraden}

\erg{26}{a) $(3 | 0 | -2)$ – der Vektor, der mit dem Parameter $t$ multipliziert wird \quad b) $g\colon \vec x = (2 | 1 | 3) + r \cdot (2 | 3 | -1)$}
\erg{27}{a) $\vec x = (2 | -4 | 1) + \lambda \cdot (-8 | 24 | 12)$ mit $0 \le \lambda \le 1$: die Liftstrecke von der Talstation ($\lambda = 0$) bis zur Bergstation ($\lambda = 1$), nicht die ganze Gerade. \quad b) Die Rutschstrecke beginnt am Start $(0 | 5 | 12)$ in 12 m Höhe ($\lambda = 0$) und endet am Ziel $(15 | -5 | 2)$ in 2 m Höhe ($\lambda = 1$); der Vektor $(15 | -10 | -10)$ führt vom Start zum Ziel. \quad c) Das Band ist eine Strecke: es beginnt bei $(3 | 1 | 0)$ am Boden ($\mu = 0$) und endet bei $(9 | 4 | 6)$ in 6 m Höhe ($\mu = 1$); $(6 | 3 | 6)$ ist der Weg vom Anfang zum Ende. \quad d) zum Beispiel $p\colon \vec x = (4 | 1 | 2) + r \cdot (1 | 3 | -2)$ – der Stützpunkt liegt nicht auf $g$, weil $(1 | 0 | 0)$ kein Vielfaches von $(1 | 3 | -2)$ ist; $q\colon \vec x = (3 | 1 | 2) + r \cdot (2 | 0 | 1)$ – Stützpunkt von $g$ und $2 + 0 - 2 = 0$. \quad e) zum Beispiel $p\colon \vec x = (-2 | 4 | 3) + t \cdot (2 | -1 | 2)$ – der Stützpunkt liegt nicht auf $g$, weil $(0 | 0 | 2)$ kein Vielfaches von $(2 | -1 | 2)$ ist; $q\colon \vec x = (-2 | 4 | 1) + t \cdot (1 | 2 | 0)$ – Stützpunkt von $g$ und $2 - 2 + 0 = 0$. \quad f) zum Beispiel $p\colon \vec x = (1 | -3 | 6) + r \cdot (3 | 1 | -1)$ – der Stützpunkt liegt nicht auf $g$, weil $(0 | 0 | 1)$ kein Vielfaches von $(3 | 1 | -1)$ ist; $q\colon \vec x = (1 | -3 | 5) + r \cdot (1 | 0 | 3)$ – Stützpunkt von $g$ und $3 + 0 - 3 = 0$. \quad g) $T = P + \tfrac{1}{4} \cdot \overrightarrow{PQ} = (4 | 0 | 2)$, also $\vec x = (4 | 0 | 2) + r \cdot (1 | 4 | 2)$ \quad h) $T = P + \tfrac{1}{3} \cdot \overrightarrow{PQ} = (-1 | 2 | 3)$, also $\vec x = (-1 | 2 | 3) + t \cdot (2 | -1 | 1)$ \quad i) $T = P + \tfrac{3}{4} \cdot \overrightarrow{PQ} = (6 | 1 | 4)$, also $\vec x = (6 | 1 | 4) + r \cdot (1 | 1 | -2)$ \quad j) Richtungsvektor $1 \cdot (1 | 0 | 2) + 2 \cdot (0 | 3 | -1) = (1 | 6 | 0)$ mit dritter Komponente null, also $g\colon \vec x = (3 | 2 | 2) + \lambda \cdot (1 | 6 | 0)$ \quad k) Richtungsvektor $2 \cdot (2 | 1 | -1) + 1 \cdot (1 | -1 | 2) = (5 | 1 | 0)$ mit dritter Komponente null, also $g\colon \vec x = (4 | 1 | 6) + \lambda \cdot (5 | 1 | 0)$}
\erg{28}{a) aus der zweiten oder der dritten – die erste Komponente des Richtungsvektors ist null \quad b) Aus der ersten Koordinate $r = 3$; zweite: $2 + 3 = 5$ ✓, dritte: $-1 + 9 = 8$ ✓ – $P$ liegt auf $g$.}
\erg{29}{a) Jeder Punkt von $g$ hat die zweite Koordinate $-1$ (Richtungskomponente null), $A$ hat 3. \quad b) Jeder Punkt von $h$ hat die dritte Koordinate 6 (Richtungskomponente null), $B$ hat 5. \quad c) Jeder Punkt von $k$ hat die erste Koordinate $-3$ (Richtungskomponente null), $C$ hat 1 – dass die zweite und die dritte Koordinate zu $t = 3$ passen, ändert daran nichts. \quad d) $1 - r = -2$ liefert $r = 3$; $y_B = 3 + 6 = 9$, $z_B = 2 - 9 = -7$ \quad e) $-1 + t = 2$ liefert $t = 3$; $x_P = 0 + 6 = 6$, $z_P = 4 - 6 = -2$ \quad f) $0 + t = -3$ liefert $t = -3$; $x_Q = 3 + 6 = 9$, $y_Q = -2 - 12 = -14$ \quad g) $|\vec u| = 3$, also $A + 3 \cdot \vec u = (9 | -2 | 8)$ (oder $A - 3 \cdot \vec u = (-3 | 4 | -4)$) \quad h) $|\vec u| = 5$, also $A + 2 \cdot \vec u = (8 | 4 | -7)$ (oder $A - 2 \cdot \vec u = (-8 | 4 | 5)$) \quad i) $|\vec u| = 3$, Faktor $4{,}5 : 3 = 1{,}5$, also $A + 1{,}5 \cdot \vec u = (-0{,}5 | 4 | 6)$ \quad j) Ansatz $(2 | 6 | 0) = (8 | 0 | 0) + t \cdot (-8 | 8 | 0)$: erste und zweite Koordinate liefern $t = 0{,}75$, die dritte stimmt, und $0 \le 0{,}75 \le 1$ – $T$ liegt auf der Strecke. \quad k) Ansatz $(-3 | 7 | 0) = (6 | -2 | 0) + t \cdot (-4 | 4 | 0)$: erste und zweite Koordinate liefern $t = 2{,}25$, die dritte stimmt – $T$ liegt auf der Geraden $EF$, wegen $2{,}25 > 1$ aber nicht auf dem Zaun. \quad l) Ansatz $L = A + t \cdot (6 | 3 | -3)$: aus der ersten Koordinate $t = \tfrac{2}{3}$; zweite: $1 + 2 = 3$ ✓, dritte: $4 - 2 = 2$ ✓; $0 \le \tfrac{2}{3} \le 1$ – $L$ liegt auf dem Seil.}
\erg{30}{a) $\overrightarrow{AB} = (2 | 1 | -2)$, $\overrightarrow{AC} = (6 | 3 | -6) = 3 \cdot \overrightarrow{AB}$ – die Verbindungsvektoren sind kollinear, $A$, $B$, $C$ liegen auf einer Geraden. \quad b) $\overrightarrow{AB} = (1 | -2 | 1)$, $\overrightarrow{AC} = (-2 | 4 | -2) = -2 \cdot \overrightarrow{AB}$ – kollinear, die drei Punkte liegen auf einer Geraden. \quad c) $\overrightarrow{AC} = (6 | 3 | 10)$, aber $3 \cdot \overrightarrow{AB} = (6 | 3 | 9)$ – kein Vielfaches, die Punkte liegen nicht auf einer Geraden. \quad d) $|\overrightarrow{AE}| = 21$ LE, 84 m sind $16{,}8$ LE, also $\lambda = 16{,}8 : 21 = 0{,}8$; $x_3 = 2 + 0{,}8 \cdot 9 = 9{,}2$ LE, Höhe: $9{,}2 \cdot 5 = 46$ m \quad e) $|\overrightarrow{AE}| = 12$ LE, 30 m sind 3 LE, also $\lambda = 0{,}25$; $x_3 = 4 + 0{,}25 \cdot 8 = 6$ LE, Höhe: $6 \cdot 10 = 60$ m \quad f) $|\overrightarrow{AE}| = 12$ LE, 180 m sind 9 LE, also $\lambda = 0{,}75$; $x_3 = 0 + 0{,}75 \cdot 8 = 6$ LE, Höhe: $6 \cdot 20 = 120$ m \quad g) $\vec x = (5 | 0 | 0) + r \cdot (-4 | 4 | 6)$; aus der ersten Koordinate $r = 0{,}5$, zweite: $0 + 2 = 2$ ✓, dritte: $0 + 3 = 3 \ne 5$ – $S$ liegt nicht auf der Geraden. \quad h) $\vec x = (0 | 6 | 1) + r \cdot (4 | -6 | 3)$; aus der ersten Koordinate $r = 0{,}5$, zweite: $6 - 3 = 3$ ✓, dritte: $1 + 1{,}5 = 2{,}5 \ne 2$ – $S$ liegt nicht auf der Geraden. \quad i) $\vec x = (2 | 0 | 0) + r \cdot (-4 | 6 | 6)$; aus der ersten Koordinate $r = 0{,}5$, zweite: $0 + 3 = 3$ ✓, dritte: $0 + 3 = 3$ ✓ – $S$ liegt auf der Geraden (in der Mitte zwischen $A$ und $G$).}
\erg{31}{a) $P_t = (t | 1 + 2t | 3 + 2t)$; $|\overrightarrow{FP_t}|^2 = (t - 4)^2 + (2t + 4)^2 + (2t - 2)^2 = 9t^2 + 36 = 56{,}25$ liefert $t = 1{,}5$ oder $t = -1{,}5$: $(1{,}5 | 4 | 6)$ und $(-1{,}5 | -2 | 0)$}
\erg{32}{a) ob ein Stützpunkt von $g$ auf $h$ liegt – er entscheidet zwischen identisch und echt parallel \quad b) $(3 | 1 | 0)$ ist kein Vielfaches von $(1 | 3 | 0)$: die erste Koordinate verlangt den Faktor 3, die zweite den Faktor $\tfrac{1}{3}$ – verschiedene Richtungen; $g$ und $h$ schneiden sich im gemeinsamen Stützpunkt und sind nicht identisch. \quad c) $\overrightarrow{AB} = (30 | 20 | 3)$, $\overrightarrow{CD} = (60 | 40 | 6) = 2 \cdot \overrightarrow{AB}$ – die Tunnelachsen sind parallel.}
\erg{33}{a) $\overrightarrow{PQ} = (4 | -2 | 5)$ ist kein Vielfaches von $(2 | -1 | 3)$ (Faktor 2 in den ersten beiden Koordinaten, $\tfrac{5}{3}$ in der dritten) – $Q$ liegt nicht auf $g$, die Geraden sind echt parallel und nicht identisch. \quad b) $\overrightarrow{PQ} = (2 | 1 | 4)$ ist kein Vielfaches von $(1 | 0 | 2)$ (die zweite Koordinate ist 1, beim Richtungsvektor 0) – $Q$ liegt nicht auf $g$, die Geraden sind echt parallel und nicht identisch. \quad c) $\overrightarrow{PQ} = (6 | -3 | 5)$ ist kein Vielfaches von $(-2 | 1 | -2)$ (Faktor $-3$ in den ersten beiden Koordinaten, $-2{,}5$ in der dritten) – $Q$ liegt nicht auf $g$, die Geraden sind echt parallel und nicht identisch. \quad d) $h\colon \vec x = (5 | 5 | 0) + t \cdot (-5 | 0 | 5)$; $AC$ liegt in der $x_1x_2$-Ebene, $h$ durchstößt diese Ebene nur in $B$, und $B$ liegt nicht auf $AC$; die Richtungen $(-5 | 5 | 0)$ und $(-5 | 0 | 5)$ sind nicht parallel – windschief. \quad e) $h\colon \vec x = (8 | 8 | 0) + t \cdot (-4 | -4 | 6)$; $BD$ liegt in der $x_1x_2$-Ebene, $h$ durchstößt diese Ebene nur in $C$, und $C$ liegt nicht auf $BD$; die Richtungen $(-8 | 8 | 0)$ und $(-4 | -4 | 6)$ sind nicht parallel – windschief. \quad f) $h\colon \vec x = (0 | 5 | 0) + t \cdot (3 | -5 | 3)$; $PT$ liegt in der $x_1x_3$-Ebene, $h$ durchstößt diese Ebene nur in $U$, und $U$ liegt nicht auf $PT$; die Richtungen $(-3 | 0 | 3)$ und $(3 | -5 | 3)$ sind nicht parallel – windschief. \quad g) Beide Geraden liegen in $E$; die Richtungen $(2 | 2 | 4)$ und $(2 | 2 | -2)$ sind keine Vielfachen – zwei nicht parallele Geraden einer Ebene schneiden sich. \quad h) Alle vier Punkte haben die dritte Koordinate 2, beide Geraden liegen also in der Ebene $x_3 = 2$; $(4 | 2 | 0)$ und $(2 | -6 | 0)$ sind keine Vielfachen – die Geraden schneiden sich. \quad i) Beide Geraden liegen in $E$; die Richtungen $(3 | 4 | -6)$ und $(1 | -6 | -2)$ sind keine Vielfachen – zwei nicht parallele Geraden einer Ebene schneiden sich. \quad j) Nein. $g$ und $p$ liegen in einer Ebene, eine senkrecht schneidende Gerade muss nicht in ihr liegen: $q\colon \vec x = (1 | 3 | 1) + t \cdot (1 | 0 | -2)$ schneidet $g$ senkrecht ($2 + 0 - 2 = 0$), hat aber stets die zweite Koordinate 3, $p$ stets 5 – $q$ und $p$ sind windschief. \quad k) Nein. $q\colon \vec x = (0 | 2 | 3) + t \cdot (1 | -1 | 0)$ schneidet $g$ senkrecht ($1 - 1 + 0 = 0$), bleibt aber in der Höhe $x_3 = 3$, während $p$ in der Höhe 5 verläuft – $q$ und $p$ sind windschief.}
\erg{34}{a) Strecke – das Seil beginnt am ersten und endet am zweiten Mast \quad b) $\vec x = (0 | 0 | 8) + t \cdot (6 | 8 | -3)$ mit $0 \le t \le 1$}
\erg{35}{a) Unteres Seil: $\vec x = (0 | 0 | 6) + s \cdot (4 | 8 | -2)$. Gleiche zweite Koordinate: $8s = 2$ liefert $s = 0{,}25$; gleiche erste: $-2 + 3t = 1$ liefert $t = 1$; Punkte $(1 | 2 | 5{,}5)$ und $(1 | 2 | 10)$, Abstand $10 - 5{,}5 = 4{,}5$ m \quad b) Gleiche zweite Koordinate: $6s = 3$ liefert $s = 0{,}5$; gleiche erste: $2t = 3$ liefert $t = 1{,}5$; Punkte $(3 | 3 | 3{,}5)$ und $(3 | 3 | 6{,}5)$, Abstand $6{,}5 - 3{,}5 = 3$ m \quad c) Gleiche erste und zweite Koordinate: $1 + 2s = t$ und $1 + 2s = 6 - t$ liefern $t = 3$, $s = 1$; Punkte $(3 | 3 | 2)$ und $(3 | 3 | 6)$, Abstand $6 - 2 = 4$ m \quad d) Horizontalabstand $\sqrt{36 + 64} = 10$ m, Höhenunterschied $1{,}5$ m, Neigung $\frac{1{,}5}{10} = 15\,\%$ \quad e) Horizontalabstand $\sqrt{144 + 25} = 13$ m, Höhenunterschied $2{,}6$ m, Neigung $\frac{2{,}6}{13} = 20\,\%$ \quad f) Horizontalabstand $\sqrt{64 + 225} = 17$ m, Höhenunterschied $1{,}7$ m, Neigung $\frac{1{,}7}{17} = 10\,\%$ \quad g) Aus dem bekannten Teilverhältnis den Auflagepunkt $X$ auf $UV$ bestimmen; die Seilgerade durch den Befestigungspunkt am ersten Mast und $X$ aufstellen; den zweiten Mast als senkrechte Gerade durch $M$ beschreiben (erste und zweite Koordinate wie bei $M$) und mit der Seilgeraden schneiden; der gesuchte Abstand ist die dritte Koordinate des Schnittpunkts minus die Dachhöhe – eine Höhendifferenz, kein Lotabstand zu einer Ebene. \quad h) Den Auflagepunkt $X$ aus dem Teilverhältnis auf $KL$ bestimmen; die Seilgerade durch $S$ und $X$ aufstellen; den Pfosten als senkrechte Gerade durch $P$ (erste und zweite Koordinate wie bei $P$) mit der Seilgeraden schneiden; von der dritten Koordinate des Schnittpunkts die Mauerhöhe abziehen. \quad i) $K_t$ ist die Position der Kamera auf der Schiene; der Parameter $t$ zwischen null und eins hält sie zwischen $A$ und $B$. Mit $r$ geht man von dieser Position aus in Blickrichtung – das ist die Blickgerade. Die Gleichung fragt, ob $F$ für eine Position auf der Schiene auf der Blickgeraden liegt; hat sie eine Lösung mit $t$ zwischen null und eins und positivem $r$, kann die Kamera $F$ filmen. \quad j) Der Teil mit $t$ beschreibt die Position des Scheinwerfers auf der Schiene von $P$ nach $Q$ – deshalb liegt $t$ zwischen null und eins. Mit $r$ kommt der Lichtstrahl in Strahlrichtung dazu. Die Gleichung fragt, ob es eine Position auf der Schiene gibt, von der aus der Strahl $Z$ trifft; $t$ und $r$ sind zwei verschiedene Parameter.}
\erg{36}{a) Horizontalabstand $\frac{0{,}9 - 0{,}3}{0{,}04} = 15$ m; linkes Ende bei $x = -15$, also $15 - 10 = 5$ m vom Rand entfernt \quad b) Horizontalabstand $0{,}45 : 0{,}06 = 7{,}5$ m; die Rampe ragt $7{,}5 - 4 = 3{,}5$ m in den Gehweg hinein \quad c) Horizontalabstand $\frac{12 - 3}{0{,}05} = 180$ m; Endpunkt $(180 | 0 | 3)$}
